\chapter{Gegnerkästen}

Um einen Gegnerkasten zu bauen kann die Umgebung \texttt{gegner} verwendet werden.

\section{Syntax}

\begin{verbatim}
\begin{gegner}
    \name{Name des Gegners oder Gegnertyps}

    \MU{Mut-Wert}
    \KL{Klugheits-Wert}
    % Dies wiederholt sich für alle Attribute. Jedes Attribut ist optional

    \waffe{Waffenname}
    \AT{AT-Wert der Waffe}
    \PA{Optional: PA-Wert der Waffe}
    \TP{Schaden der Waffe}
    \INI{Initiative}

    \ruestung{Name der Rüstung}
    \RS{Optional: RS der Rüstung}
\end{gegner}
\end{verbatim}

\section{Beispiel}

\subsection{Quellcode}

\begin{verbatim}
\begin{gegner}
    \name{Soldat}
    \KK{12}

    \waffe{Schwert}
    \AT{12}
    \PA{4}
    \TP{1W+1}
    \INI{12}

    \ruestung{Lederharnisch}
    \RS{2}
\end{gegner}
\end{verbatim}

\subsection{Ausgabe}

\begin{gegner}
    \name{Soldat}
    \KK{12}

    \waffe{Schwert}
    \AT{12}
    \PA{4}
    \TP{1W+1}
    \INI{12}

    \ruestung{Lederharnisch}
    \RS{2}
\end{gegner}